\documentclass[12pt]{article}
\usepackage[a4paper,margin=1.25cm]{geometry}
\usepackage{amsmath,amssymb,mathtools}
\usepackage{graphicx}
\usepackage{booktabs}
\usepackage{tabularx}
\usepackage{tikz}
\usepackage[T1]{fontenc}
\usepackage[utf8]{inputenc}
\usepackage[english]{babel}

\setlength{\parindent}{0pt}
\setlength{\parskip}{0.45em}
\setlength{\abovedisplayskip}{5pt plus 1pt minus 1pt}
\setlength{\belowdisplayskip}{5pt plus 1pt minus 1pt}
\setlength{\abovedisplayshortskip}{3pt plus 1pt minus 1pt}
\setlength{\belowdisplayshortskip}{3pt plus 1pt minus 1pt}

\newcommand{\dd}{\,\mathrm{d}}
\newcommand{\ii}{\mathrm{i}}
\newcommand{\e}{\mathrm{e}}
\newcommand{\om}{\omega}

\begin{document}

\begin{center}
{\Large\bfseries A logarithm over the cyclotomic quadratic}\\[0.35em]
{\large one parameter, one partial fraction, and a rational function that collapses}
\end{center}

\[
\boxed{\displaystyle
\int_0^\infty\frac{\log(1+x^3)}{1+x+x^2}\dd x
=\frac{2\pi\log2}{\sqrt3}.}
\]

\textbf{Overview.}
Introduce the parameter
\[
F(t)=\int_0^\infty\frac{\log(1+t^3x^3)}{1+x+x^2}\dd x,
\qquad t\ge0,
\]
so that the problem is \(F(1)\). Differentiating in \(t\) turns the
logarithm into a rational function, and that rational integral has a
simple value: the whole cubic--quadratic interaction cancels
and leaves
\[
\int_0^\infty\frac{\dd x}{(1+t^3x^3)(1+x+x^2)}
=\frac{2\pi}{3\sqrt3}\cdot\frac1{1+t}.
\]
Hence \(F'(t)=\frac{2\pi}{\sqrt3}\cdot\frac1{1+t}\), and integrating back
from \(F(0)=0\) gives
\[
\boxed{\displaystyle
F(t)=\frac{2\pi}{\sqrt3}\log(1+t),\qquad t\ge0,}
\]
of which the stated integral is the case \(t=1\).

The reason the constant \(\frac{2\pi}{3\sqrt3}\) governs everything is
that \(\frac{\dd x}{1+x+x^2}\) is, up to the factor \(\frac1{\sqrt3}\),
arclength on a circular arc; Section 5 makes this precise.

\section*{1. The parametrised integral}

For \(t>0\) put
\[
F(t)=\int_0^\infty\frac{\log(1+t^3x^3)}{1+x+x^2}\dd x .
\]
Convergence is immediate: at \(0\) the integrand vanishes like
\(t^3x^3\), and at infinity it behaves like \(\frac{3\log x}{x^2}\).
Also \(F(0)=0\).

Differentiating under the integral sign,
\[
F'(t)=\int_0^\infty\frac{3t^2x^3}{(1+t^3x^3)(1+x+x^2)}\dd x .
\tag{1}
\]
The differentiation is legitimate on any interval \(0<t_1\le t\le t_2\):
the \(t\)-derivative of the integrand is bounded by
\(\frac{3}{t_1}\cdot\frac1{1+x+x^2}\), because
\(\frac{t^3x^3}{1+t^3x^3}\le1\), and that bound is integrable.

The identity
\[
3t^2x^3=\frac{3}{t}\bigl[(1+t^3x^3)-1\bigr]
\]
splits (1) into
\[
F'(t)=\frac3t\left[
\int_0^\infty\frac{\dd x}{1+x+x^2}
-\int_0^\infty\frac{\dd x}{(1+t^3x^3)(1+x+x^2)}
\right]
=\frac3t\bigl[C-K(t)\bigr],
\tag{2}
\]
where
\[
C=\int_0^\infty\frac{\dd x}{1+x+x^2},
\qquad
K(t)=\int_0^\infty\frac{\dd x}{(1+t^3x^3)(1+x+x^2)} .
\]

\section*{2. The base constant}

Completing the square,
\[
C=\int_0^\infty\frac{\dd x}{\left(x+\frac12\right)^2+\frac34}
=\frac{2}{\sqrt3}\left[\arctan\frac{2x+1}{\sqrt3}\right]_0^\infty
=\frac{2}{\sqrt3}\left(\frac\pi2-\frac\pi6\right),
\]
so
\[
\boxed{\displaystyle C=\int_0^\infty\frac{\dd x}{1+x+x^2}
=\frac{2\pi}{3\sqrt3}.}
\tag{3}
\]

\section*{3. The rational lemma}

\textbf{Lemma.} For every \(t>0\),
\[
\boxed{\displaystyle
K(t)=\int_0^\infty\frac{\dd x}{(1+t^3x^3)(1+x+x^2)}
=\frac{2\pi}{3\sqrt3}\cdot\frac1{1+t}
=\frac{C}{1+t}.}
\tag{4}
\]

\textit{Proof.} Write \(P(x)=1+t^3x^3\) and \(Q(x)=1+x+x^2\), and seek the
partial fraction decomposition
\[
\frac1{P(x)Q(x)}=\frac{\alpha(x)}{P(x)}+\frac{\beta(x)}{Q(x)},
\qquad \deg\alpha\le2,\ \deg\beta\le1 .
\]
Equivalently \(1=\alpha Q+\beta P\). Evaluate at the roots of \(Q\),
namely the primitive cube roots of unity \(\om=\e^{2\pi\ii/3}\) and
\(\overline\om\). Since \(\om^3=\overline\om^{\,3}=1\),
\[
P(\om)=1+t^3\om^3=1+t^3=P(\overline\om),
\]
so
\[
\beta(\om)=\beta(\overline\om)=\frac1{1+t^3}.
\]
A polynomial of degree \(\le1\) taking equal values at two distinct
points is constant, hence
\[
\beta(x)\equiv\frac1{1+t^3}.
\tag{5}
\]
This is the only place where the cubic structure is used, and it is
exactly the statement that \(Q\) divides \(x^3-1\).

For \(\alpha\), subtract:
\[
\frac{\alpha(x)}{P}
=\frac1{PQ}-\frac1{(1+t^3)Q}
=\frac{(1+t^3)-P(x)}{(1+t^3)PQ}
=\frac{t^3\bigl(1-x^3\bigr)}{(1+t^3)PQ}.
\]
Now \(1-x^3=(1-x)(1+x+x^2)=(1-x)Q\), so the factor \(Q\) cancels and
\[
\alpha(x)=\frac{t^3(1-x)}{1+t^3}.
\tag{6}
\]
Thus
\[
\frac1{P(x)Q(x)}
=\frac1{1+t^3}\left[\frac{t^3(1-x)}{1+t^3x^3}+\frac1{1+x+x^2}\right].
\tag{7}
\]
Both pieces are separately integrable on \([0,\infty)\): the first decays
like \(-\frac1{t^3x^2}\).

For the first piece substitute \(u=tx\):
\[
\int_0^\infty\frac{1-x}{1+t^3x^3}\dd x
=\frac1t\int_0^\infty\frac{\dd u}{1+u^3}
-\frac1{t^2}\int_0^\infty\frac{u\dd u}{1+u^3}.
\]
Both integrals are instances of the beta-type formula
\[
\int_0^\infty\frac{u^{s-1}}{1+u^3}\dd u=\frac{\pi/3}{\sin\frac{\pi s}{3}},
\qquad 0<s<3,
\tag{8}
\]
and \(\sin\frac\pi3=\sin\frac{2\pi}3=\frac{\sqrt3}{2}\), so both equal
\(\frac{2\pi}{3\sqrt3}=C\). (Formula (8) follows from the standard
\(\int_0^\infty\frac{u^{s-1}}{1+u^n}\dd u=\frac{\pi/n}{\sin(\pi s/n)}\);
the coincidence of the two values here is the symmetry \(s\mapsto3-s\).)
Hence
\[
\int_0^\infty\frac{1-x}{1+t^3x^3}\dd x
=C\left(\frac1t-\frac1{t^2}\right)
=C\cdot\frac{t-1}{t^2}.
\]
Substituting this and (3) into (7),
\[
K(t)=\frac{1}{1+t^3}\left[t^3\cdot C\cdot\frac{t-1}{t^2}+C\right]
=\frac{C\bigl[t(t-1)+1\bigr]}{1+t^3}
=\frac{C\,(t^2-t+1)}{(1+t)(t^2-t+1)}
=\frac{C}{1+t},
\]
using \(1+t^3=(1+t)(t^2-t+1)\). \(\square\)

Everything turns on the cancellation of \(t^2-t+1\). The numerator
\(t(t-1)+1\) produced by the two beta integrals is literally the same
cyclotomic factor that appears in the factorisation of \(1+t^3\); nothing
survives except the linear factor \(1+t\), which is what makes the
antiderivative a logarithm.

\section*{4. Conclusion}

Insert (3) and (4) into (2):
\[
F'(t)=\frac3t\left[C-\frac{C}{1+t}\right]
=\frac{3C}{t}\cdot\frac{t}{1+t}
=\frac{3C}{1+t}
=\frac{2\pi}{\sqrt3}\cdot\frac1{1+t}.
\tag{9}
\]
Since \(F(0)=0\),
\[
F(t)=\int_0^t\frac{2\pi}{\sqrt3}\frac{\dd s}{1+s}
=\frac{2\pi}{\sqrt3}\log(1+t),
\]
that is,
\[
\boxed{\displaystyle
\int_0^\infty\frac{\log(1+t^3x^3)}{1+x+x^2}\dd x
=\frac{2\pi}{\sqrt3}\log(1+t),\qquad t\ge0 .}
\tag{10}
\]
Setting \(t=1\),
\[
\int_0^\infty\frac{\log(1+x^3)}{1+x+x^2}\dd x=\frac{2\pi\log2}{\sqrt3}
=2.5144598308\ldots,
\]
in agreement with numerical quadrature to the displayed precision.

Two consistency checks on (10). As \(t\to0^+\), the left side is
\(\approx t^3\int_0^\infty\frac{x^3}{1+x+x^2}\dd x\), which diverges,
while the right side is \(\approx\frac{2\pi t}{\sqrt3}\). There is no
contradiction, because for small \(t\) the logarithm is only
\emph{eventually} small: the crossover happens at \(x\sim1/t\), and it is
the tail beyond it that produces the linear behaviour. As
\(t\to\infty\), \(\log(1+t^3x^3)\approx3\log t+3\log x\) and
\(3C\log t=\frac{2\pi}{\sqrt3}\log t\) matches the leading term of
\(\frac{2\pi}{\sqrt3}\log(1+t)\).

\section*{5. Why \(2\pi/(3\sqrt3)\): the arc behind the measure}

The constant \(C\) is not arbitrary. Let \(\om=\e^{2\pi\ii/3}\), so that
\(1+x+x^2=(x-\om)(x-\overline\om)\), and set
\[
z=\frac{x-\om}{x-\overline\om}.
\]
For real \(x\) the two distances \(|x-\om|\) and \(|x-\overline\om|\) are
equal, so \(|z|=1\): the real line is mapped to the unit circle. Moreover
\[
\frac{\dd z}{z}
=\left(\frac1{x-\om}-\frac1{x-\overline\om}\right)\dd x
=\frac{\overline\om-\om}{(x-\om)(x-\overline\om)}\dd x
=\frac{-\ii\sqrt3}{1+x+x^2}\dd x,
\]
since \(\om-\overline\om=\ii\sqrt3\). Writing \(z=\e^{\ii\varphi}\),
\(\frac{\dd z}{z}=\ii\dd\varphi\), so
\[
\boxed{\displaystyle
\frac{\dd x}{1+x+x^2}=-\frac{\dd\varphi}{\sqrt3}.}
\tag{11}
\]
The measure in the problem \emph{is} arclength on the unit circle, scaled
by \(\frac1{\sqrt3}\).

\begin{center}
\begin{tikzpicture}[scale=2.6,>=stealth]
\draw[->] (-1.35,0)--(1.4,0) node[right] {\scriptsize \(\operatorname{Re}z\)};
\draw[->] (0,-1.3)--(0,1.35) node[above] {\scriptsize \(\operatorname{Im}z\)};
\draw[dashed] (0,0) circle (1);
\draw[very thick,domain=-120:0,samples=60] plot ({cos(\x)},{sin(\x)});
\fill (1,0) circle (0.9pt) node[right] {\scriptsize \(z=1\ (x=\infty)\)};
\fill ({cos(-120)},{sin(-120)}) circle (0.9pt)
   node[below left] {\scriptsize \(z=\overline\om\ (x=0)\)};
\fill ({cos(-60)},{sin(-60)}) circle (0.9pt)
   node[below right] {\scriptsize \(x=1\)};
\draw[->] (0.42,0) arc (0:-120:0.42);
\node at (0.16,-0.52) {\scriptsize \(\tfrac{2\pi}{3}\)};
\node at (0,-1.5) {\scriptsize total arc \(\tfrac{2\pi}{3}\), hence \(C=\tfrac1{\sqrt3}\cdot\tfrac{2\pi}{3}\)};
\end{tikzpicture}
\end{center}

As \(x\) runs from \(0\) to \(\infty\), \(z\) runs from
\(\frac{-\om}{-\overline\om}=\om^{2}=\overline\om\) to \(1\), covering an
arc of angular length \(\frac{2\pi}{3}\), exactly one third of the
circle, because \(\om\) is a primitive cube root of unity. Therefore
\[
C=\frac1{\sqrt3}\cdot\frac{2\pi}{3}=\frac{2\pi}{3\sqrt3},
\]
recovering (3) with no antiderivative at all. The same picture explains
why the answer to the whole problem carries a factor \(\frac{\pi}{\sqrt3}\)
rather than, say, \(\pi\): the \(\sqrt3\) is \(|\om-\overline\om|\), and
the \(\frac{2\pi}{3}\) is the third of the circle cut out by the
cyclotomic quadratic.

\section*{6. The pattern for other exponents}

Nothing forces the exponent \(3\). The same three steps, writing
\(nt^{n-1}x^n=\frac{n}{t}\bigl[(1+t^nx^n)-1\bigr]\), decomposing, and using
(8), apply whenever the denominator \(Q\) divides \(x^n-1\), since
that is what made \(\beta\) constant in (5).

\begin{center}
\small
\renewcommand{\arraystretch}{1.45}
\begin{tabularx}{\textwidth}{@{}>{\raggedright\arraybackslash}p{0.24\textwidth}>{\raggedright\arraybackslash}p{0.24\textwidth}X@{}}
\toprule
\(Q(x)\) & divides & resulting identity \\
\midrule
\(1+x\) & \(x^2-1\) &
\(\displaystyle\int_0^\infty\frac{\log(1+t^2x^2)}{1+x}\dd x\)
diverges; the quadratic \(Q\) is needed for convergence. \\
\addlinespace
\(1+x+x^2\) & \(x^3-1\) &
\(\displaystyle\int_0^\infty\frac{\log(1+t^3x^3)}{1+x+x^2}\dd x
=\frac{2\pi}{\sqrt3}\log(1+t)\). \\
\addlinespace
\(1+x^2\) & \(x^4-1\) &
\(\displaystyle\int_0^\infty\frac{\log(1+t^4x^4)}{1+x^2}\dd x
=\pi\log(1+\sqrt2\,t+t^2)\), by the same route with
\(\int_0^\infty\frac{u^{s-1}}{1+u^4}\dd u=\frac{\pi/4}{\sin(\pi s/4)}\). \\
\bottomrule
\end{tabularx}
\end{center}

In each case the mechanism is identical: the roots of \(Q\) are roots of
unity, so \(P\) takes the same value \(1+t^n\) at all of them, the
correction term \(\alpha\) inherits the factor \(Q\) from
\(1-x^n\), and the beta integrals supply a numerator that cancels the
non-linear cyclotomic factor of \(1+t^n\). What is left is always a
logarithm.

\end{document}
